\documentclass{article}


\usepackage[a4paper, margin=2cm]{geometry}
%
\usepackage{amsfonts, amsmath, amssymb}
\usepackage{mathrsfs, mathtools, mathbbol}
\usepackage{tikz}
\usepackage{tikz-cd}
%
\usepackage[T2A]{fontenc}
\usepackage[utf8]{inputenc}
\usepackage[english]{babel}
\usepackage{alphabeta}
\usepackage{hyperref}

\usepackage[backend=biber]{biblatex}
\usetikzlibrary{shapes,arrows, positioning,fit}

\newcommand*{\sheaf}{\mathscr}


\usepackage[most]{tcolorbox}

\newtcbtheorem[number within=section]{definition}{Definition}{
    enhanced,
    breakable,
    colback=gray!5,
    colframe=gray!50,
    fonttitle=\bfseries,
    separator sign={.},
    before skip=10pt,
    after skip=10pt,
    boxrule=0.5pt,
    arc=2pt,
    left=8pt,
    right=8pt,
    top=6pt,
    bottom=6pt
}{def}


\newtcbtheorem[number within=section]{proposition}{Proposition}{
    enhanced,
    breakable,
    colback=gray!5,
    colframe=gray!50,
    fonttitle=\bfseries,
    separator sign={.},
    before skip=10pt,
    after skip=10pt,
    boxrule=0.5pt,
    arc=2pt,
    left=8pt,
    right=8pt,
    top=6pt,
    bottom=6pt
}{prop}


\newtcbtheorem[number within=section]{example}{Example}{
    enhanced,
    breakable,
    colback=gray!5,
    colframe=gray!50,
    fonttitle=\bfseries,
    separator sign={.},
    before skip=10pt,
    after skip=10pt,
    boxrule=0.5pt,
    arc=2pt,
    left=8pt,
    right=8pt,
    top=6pt,
    bottom=6pt
}{ex}


\newtcbtheorem[number within=section]{remark}{Remark}{
    enhanced,
    breakable,
    colback=gray!5,
    colframe=gray!50,
    fonttitle=\bfseries,
    separator sign={.},
    before skip=10pt,
    after skip=10pt,
    boxrule=0.5pt,
    arc=2pt,
    left=8pt,
    right=8pt,
    top=6pt,
    bottom=6pt
}{remark}


\newtcbtheorem[number within=section]{note}{Note}{
    enhanced,
    breakable,
    colback=gray!5,
    colframe=gray!50,
    fonttitle=\bfseries,
    separator sign={.},
    before skip=10pt,
    after skip=10pt,
    boxrule=0.5pt,
    arc=2pt,
    left=8pt,
    right=8pt,
    top=6pt,
    bottom=6pt
}{note}


\newtcbtheorem[number within=section]{conjecture}{Conjecture}{
    enhanced,
    breakable,
    colback=gray!5,
    colframe=gray!50,
    fonttitle=\bfseries,
    separator sign={.},
    before skip=10pt,
    after skip=10pt,
    boxrule=0.5pt,
    arc=2pt,
    left=8pt,
    right=8pt,
    top=6pt,
    bottom=6pt
}{conjecture}



\newtcbtheorem[number within=section]{question}{Question}{
    enhanced,
    breakable,
    colback=gray!5,
    colframe=gray!50,
    fonttitle=\bfseries,
    separator sign={.},
    before skip=10pt,
    after skip=10pt,
    boxrule=0.5pt,
    arc=2pt,
    left=8pt,
    right=8pt,
    top=6pt,
    bottom=6pt
}{question}






\addbibresource{dgla-stack.bib}

\title{DG Lie Algebroid and Quotient Stack}
\author{Koji Sunami}


\begin{document}

\maketitle



In this paper, I'll discuss foliation from point of view of derived algebraic geometry.
Derived algebraic geometry proves the $\infty$-equivalence of dg Lie algebroid and foliation. 
Now the quotient is defined by quotient of dg $L_\infty$ groupoid generated from dg Lie albebroid, 
and it’s calculated by dg $L_\infty$ version of Chevalley-Eilenberg complex. 
I’m also wondering evaluation of the foliation using language of DT/GW invariants. 
First of all, the idea of foliation is classically defined differential geometrically on a smooth manifold: \\


\begin{definition}{Foliation}{foliation}
Foliation is a partition of a smooth manifold $M$ into its submanifolds.
One idea of generating foliation is to retrieve from distribution $D_x\subset TM$ for some $x\in M$,
and this distribution classfies the points $x\in M$ to define foliaton.
Distribution $D_x$ is defined by its vector field $X,Y\in \Gamma(D_x)$ satisfying $[X,Y]\in \Gamma(D_x)$,
and in this case $D_x$ is called integrable. This $D_x$ corresponds to a submanifold $L$,
since there's its restriction $D_x|_L = T_x L$ by the unique maximal submanifold $L\subset M$ s.t. $x\in L$.
Note in fact $D_x|_{L_y} = T_y L$ for all $y\in L$.
Accordingly, the resulting set $\{L_y\}$ for $y\in M$ is the set of submanifolds $L_y\subset M$ and disjoint union $\coprod L_y = M$.
\end{definition}


To verbalize the notion of foliation is an equivalent notion with integrability by Frobenius theorem,
where integrability is done by Lie algebra. If Lie algebra describes the global property (there's no "if" in math Toynbee says),
the corresponding Lie group quotient of manifold $[M/G]$ is a stack. It doesn't work, however. 
I'll use Lie algebroid to describe every point of the submanifold $L\subset M$,
whose Lie groupoid $\mathcal{G}$ quotient is a quotient stack, which describes foliation $[M/\mathcal{G}]$. 
I'll first of all terminology to define foliation derived algebraic geometrically.


\begin{definition}{Groupoid}{groupoid}
A groupoid is a category in which every morphisms are invertible. 
Denote the objects are $Ob(\mathcal{C}) = G_0$, and the morphisms are $Mor(\mathcal{C})=G_1$, the double arrow
$G_1 \rightrightarrows G_0$. For an arrow $g:x\to y$, there are source and target maps $s,t,:G_1\to G_0$, 
together with composition $G_1 \times_{G_0} G_1\to G_1$, identity arrow and inverses.
A group becomes a Lie groupoid if $G_0$ and $G_1$ are smoooth manifolds, and all the maps are smooth.
\end{definition}

There are many important examples of Lie groupoid. The trivials things 
(logically equivalent to Lie group but somehow interesting generalization) are Lie group,
which is a Lie groupoid with one object so $G\rightrightarrows *$.
Another example which is more geometric is that Lie group action $G\times X\rightrightarrows X$ can be modeled by $G_1=G\times X$ and $G_0 = X$.
By the context of foliation, Lie groupoid $Hol(\mathcal{F})$ of the foliation $\mathcal{F}$ is holonomy groupoid. 
All the above examples are important. Now we consider quotient stack by Lie groupoid, in particular by Lie group:


\begin{definition}{Quotient Stack}{quotient stack}
For an algebraic variety $X$, quotient stack by Lie group $G$ is $[X/G]$ is a functor $[X/G]:(Spaces)^{op}\to Grouoids$.
\[
[X/G](S) = \{ (P\to S, f:P\to X) |  P\to S \text{ is a principal G-bundle,} f is G-equivariant \}
\]
In particular if $X=pt$, the quotient stack is $[pt/G] = BG$, 
and $BG(pt)$ is one-object groupoid whose automorphism group is $G = Hom(*,*)$. \\

Meanwhile, quotient stack by groupoid $[U/R]$ for a groupoid $R$ is a stackification $[U/R] = St([U/_p R])$ of 
category fibered in groupoids $p: U \to R$. This works for all $c\in R$, there exists $x\in U$ s.t. $p(x)=c$,
and for all $f:x\to y$, $p(f) = id_c$.
\end{definition}

where I abbreviate what the stackification means in this introductory chapter for the moment.
In any case, the important quotient stack is holonomy groupoid quotient $[M/Hol(F)]$, which is a foliation.
But how does Lie groupoid appears, where Lie algebra solves integrability condition in manifold?
Lie algebra can be generalied to Lie algebroid by its globalization, where Lie groupoid comes.
The definition of Lie algebroid is:


\begin{definition}{Lie Algebroid}{lie algebroid}
For a smooth manifold $M$, a Lie algebroid over $M$ consists of $(A\to M,[\cdot,\cdot]_A, \rho)$
where $A\to M$ is a vector bundle, $[\cdot, \cdot]_A : \Gamma(A) \times \Gamma(A)\to \Gamma(A)$ is a Lie bracket of sections,
$\rho:A \to TM$ is an anchor map. These satisfy product rule $[X,fY]_A = f[X,Y]_A + \rho(X)(f) Y$,
and the anchor preserves bracket $\rho([X,Y]_A) = [\rho(X), \rho (Y)]$.
\end{definition}

The important remark is Lie algebroid doesn't always correspond to Lie groupoid unlike the Lie algebra/group correspondence,
since Lie groupoid always make Lie algebroid by infinitesimalization but not the other way around.


\begin{conjecture}{}{}

\end{conjecture}

%\begin{definition}{}{}
%\end{definition}

\nocite{*}
\printbibliography


\end{document}



